% Split from app_statmethods.tex 2026-10-11 (note reorganization): the N-D declarations, scalar-5D reductions, extended-fiducial fractions and unfolded-event inference, moved verbatim.
\section{Statistical methods for the multidimensional results and inference}
\label{app:statmethods-nd}

This appendix gives the conventions for covariances and comparisons beyond the 2D result: rank, ndf and finite-ensemble declarations for N-D covariances, the scalar five-dimensional diagnostic reductions and interval definitions, the extended-fiducial band and pull fractions, and inference from unfolded events and selected projections. The 2D methods and their worked examples are in App.~\ref{app:statmethods}.

\subsection{Rank, ndf and finite-ensemble declarations for N-D covariances}
\label{sec:nd-declarations}

The 2D justification in \S\ref{sec:chisq} does not transfer to N-D. An earlier 5D covariance from the standard production chain has $10\,694$ reported bins, effective positive rank $\mathbf{263}$ and infinite numerical condition number. It is distinct from $C_Z$, the five-axis covariance of the configuration-(a) central value, and was built before $C_Z$ was chosen for use. Its $4825$-bin 5D$\to$4D projection retains rank $263$ exactly. The covariance is a sum of about 45 rank-one outer products, one for each two-endpoint MAT systematic band, and of the statistical and machine-learning (ML) sample blocks; such a sum cannot span $10\,694$ directions, so the rank deficiency follows from the construction. Quoting $\chi^{2}/n_{\text{reported}}$ would understate tension by roughly $40\times$ in 5D or $18\times$ in 4D.

For an N-D covariance, $\mathrm{ndf}$ is the retained rank, not $n_{\text{reported}}$. Every N-D $\chi^{2}$ must state, where it is quoted, (i) the inverse actually used, including the pseudo-inverse \texttt{rcond} or the truncation rank; (ii) the retained rank, as $\mathrm{ndf}$; (iii) the rank-truncation scan of that covariance, which is expected to lack the smoothness that supports the 2D choice in \S\ref{sec:chisq}; (iv) the identity of the covariance, since the earlier 5D covariance, its 4D projection and the published 2D block differ in rank; and (v) for each sample block in the sum, the ensemble size $N$, the biased $1/N$ or unbiased $1/(N-1)$ normalization, the effective inverted dimension $p$ after truncation, and the finite-ensemble treatment, or that none is applied. The measurements, the derivation and the requirement that N-D difference statistics use a rank cut fixed in advance are in Ref.~\cite{NoteRecords2026} (\recentry{rank-inversion}).
% The records state the cut as a requirement and give no value (docs/orchestration/RANK-AND-INVERSION-20260810.md sec. 5, "owner
% boundary"; docs/STATVAL_REPAIR.md, "PREDECLARED rank/pseudoinverse cut"), so it is stated as that requirement.

\paragraph{Finite-ensemble bias of the inverse.} Item (v) records the ensemble size of each block; without it, an uncorrected finite-ensemble bias could enter unnoticed. The inverse of a noisy sample covariance is biased \emph{upward}. For an unbiased $1/(N-1)$ estimator over $N$ independent Gaussian realizations, the standard debiasing factor is $(N-p-2)/(N-1)$. For this analysis's production convention, universe-mean centered with biased $1/N$ normalization, it is $(N-p-2)/N$. The superseded unified-throw 5D covariance has $N=160$, truncated $p=100$, giving $58/160\simeq0.363$: uncorrected precision is too large by $\simeq2.76$. For fixed residuals this inflates $\chi^{2}$ and overstates tension; confidence regions from the same over-tight covariance become over-optimistic.
% (\texttt{nd-unfolding/uq\_math.py}, ``universe-mean centered, biased 1/N'')

No generic correction is prescribed. A sum of deterministic rank-one band terms and statistical and ML sample blocks, such as the earlier 5D covariance from the standard production chain, has no single debiasing factor. The standard factor also assumes independent Gaussian realizations and a data-independent truncation dimension, neither established for a data-dependent cut. The appropriate treatment depends on the covariance and inverse actually used, and is to be determined from the quantities recorded under item (v).
% The cheap half is unconditional: record $N$ beside each block in the construction receipts.

\subsection{Scalar-5D diagnostic reductions and interval definitions}
\label{sec:gbdt-diagnostic-definitions}

The synthesis in \S\ref{sec:gbdt-model-dependence} reuses saved outputs; it involves no new unfolding, training or pseudo-experiments. Fine-grid cross-section densities are converted to integrals by multiplying by all five bin widths before aggregation. The saved trace functionals on the EW cells are differential in $(\Eavail,W)$ and are multiplied by the cell areas; those on the J and H2 cells are already integrals. Direct tensor-slice integration checks the index-map sums and full-domain rate. Reporting boundaries coincide with fine edges. The J boundaries are $\pT=(0,0.55,0.85,4.5)$, $\pPar=(1.5,3.5,6,60)$, $\Eavail=(0,0.4,1.5,100)$, $q_3=(0,1.2,2,100)$ and $W=(0,1.4,2.2,100)$; H2 uses interior boundaries $(0.7,4.5,0.8,2,1.8)$, respectively, with the same outer boundaries. Energy units are GeV and momentum units GeV/$c$. A cell is supported when both its expected reconstructed signal and selected truth weight at data exposure are at least 1000. This gives 109 J, 27 H2 and 39 EW cells; the unsupported EW11, EW36 and EW37 are omitted from the supported-cell summaries only.

For experiment $i$ and cell $f$,
\begin{equation}
 r_{if}=\frac{\hat f_i-f_{{\rm true},i}}{f_{{\rm true},i}},\qquad
 \bar r_f=\frac1n\sum_i r_{if},\qquad
 s_f^2=\frac1{n-1}\sum_i(r_{if}-\bar r_f)^2.
\end{equation}
The truth of the stored pseudo-data source varies with the MC split, so $s_f$ measures relative-residual spread, not unconditional fixed-truth cross-section variance. The saved table also reports bias-retaining $\sqrt{n^{-1}\sum_i r_{if}^2}$. Pointwise 95\% mean intervals use $t_{0.975,n-1}s_f/\sqrt n$; SD intervals use $s_f\sqrt{(n-1)/\chi^2_{0.975,n-1}}$ and $s_f\sqrt{(n-1)/\chi^2_{0.025,n-1}}$. These normal-model approximations condition on finite simulation and make no simultaneous claim across cells. A deterministic noise-free trace from one simulation receives no error band; its cell-to-cell variation is not sampling uncertainty.

The development assessment includes an interval diagnostic. It uses $\sigma_f$, the sample standard deviation of 100 refitted bootstrap replicas of the nominal development pseudo-experiment (experiment 700000). Assessment results are scored against $[\hat f_i-\sigma_f,\hat f_i+\sigma_f]$ and $[\hat f_i-1.96\sigma_f,\hat f_i+1.96\sigma_f]$, without bias correction or model envelope. The table gives integer hits, both empirical fractions, exact pointwise 95\% binomial intervals and fixed half-width relative to mean truth. It conditions on estimated bootstrap widths without propagating finite-replica uncertainty. Coverage fractions are not pooled across cells to certify them. The examples shown are a populated low-energy cell (EW7) and the corner cell EW41, already known to fail; they are not freshly selected validation targets.

The rounding study keeps the fine-grid membership fixed under uniform input perturbations within half a float32 unit in the last place. With 20 base jitters, 50 unperturbed bootstrap replicas $f_b$ and their jittered partners $f'_b$, the reported quantities are
\begin{equation}
 s_{\rm base}=\operatorname{SD}(f_j),\quad
 s_{\rm boot}=\operatorname{SD}(f_b),\quad
 s_{\rm pair}=\operatorname{SD}\!\left((f'_b-f_b)/\sqrt2\right).
\end{equation}
The consistency ratio is $s_{\rm pair}/s_{\rm base}$; the estimated numerical share is $s_{\rm pair}^2/s_{\rm boot}^2$. The paired-model $\sqrt2$ factor introduces no physical noise source. Containment, as defined before the study ran, requires a median consistency ratio in $[0.7,1.4]$ and $s_{\rm boot}^2\geq s_{\rm pair}^2$ in at least 95\% of supported cells. All supported cells meet the inequality. This finite-sample diagnostic supports the construction without proving coverage or independence from every physical uncertainty.
% Note register round n1, batch n1-F: the criterion's timing in plain words (state/s5p/contract-amendment-2-stage2-studies.json,
% study_N decision_rule, FROZEN before any Stage-2 compute).

For prior vertex $v$, data sensitivity is $d_{vf}=\hat f_{{\rm data},v}/\hat f_{{\rm data},0}-1$, with envelope $h_f=\max_v|d_{vf}|$. Its denominator differs from the simulated residual's. Neither $h_f$ nor an outer product of closure residuals is given a probability model. The reductions keep the per-cell values, the identity of each input and the exclusions, and distinguish complete outputs from checkpoints that carry no metadata. From a checkpoint, only the iterations that reproduce the originally recorded values are used; later entries are excluded. The analyst cross-checked the reductions numerically.
% Reproduction: nd-unfolding/gbdt_model_dependence/{reduce.py,analyze.py,README.md}.

\subsection{Extended-fiducial statistical-band and pull fractions}
\label{sec:fps-fractions}
For the containment diagnostic in \S\ref{sec:fps}, the independent scalar statistical-covariance band contains \SI{\covFPSanalytic}{\percent} of closure bin--toy cells and exceeds the closure-toy spread. The same-ensemble pull fraction is \SI{68.9}{\percent}, with mean absolute pull $\langle|r|\rangle\approx0.79$; it remains a diagnostic of the pulls and their Gaussianity. The split-sample fractions in the published and extension regions are \SI{68.5}{\percent} and \SI{70.0}{\percent}, respectively; both are descriptive, and neither is used to certify a covariance.
% Moved from sec_fps.tex (Validation battery, item 4), organization review 2026-10-04,
% lane 3 sec. 10 finding 2 (g4 HANDED-OFF; g5 handoffs.md H4). There it followed "Separately,".

\subsection{Inference from unfolded events and selected projections}
\label{sec:unbinned-inference}
Unfolding correlates weighted events; a downstream likelihood treating them as independent can underestimate parameter uncertainties~\cite{Desai:2025correlated}. Nominal weights do not specify uncertainty. Coherent unfolding replicas could propagate it through the fit; resampling unfolded rows independently is not equivalent to resampling inputs and repeating the learned estimator. Classifier and permutation goodness-of-fit tests require null calibration for fitting and data reuse. No new validation result accompanies these requirements.

The uncertainty target matters: separate bin or projection coverage does not imply simultaneous coverage. A fixed linear projection $y=P x$ propagates covariance as $C_y=P C_x P^{\top}$, with integrated bin widths in $P$ for densities. The identity neither calibrates that covariance nor accounts for regions chosen after inspecting data.

Recent work constructs confidence intervals for selected functionals in constrained inverse problems~\cite{Stanley:2025functionals} and joint confidence regions for multiple functionals~\cite{Batlle:2025simultaneous}, which could provide independent benchmarks for specified integrals or contrasts. Their guarantees depend on the forward model, noise law, and parameter constraints and do not transfer automatically to a learned response with finite simulation and detector nuisances. A rank-deficient forward operator is also distinct from a rank-deficient sample covariance; neither truncating the latter nor validating a few projections certifies the full distribution. Such a benchmark would supplement, not replace, the scalar five-dimensional uncertainty construction that a measurement requires.

% Methods rewrite draft 4 (X5 G9): the self-contained familywise-error derivation. Sources: docs/orchestration/state/2d-closeout/stageF/methods/M2-FWER-RECORD.md:39-74 (bound, certified table, assumptions A1/A2, the rule, replacement text); M2-AUDIT.md:76-117 (Q2: model, crossing-event inclusion, first-true-null Holm reduction without independence, claim variants, ordering); m2-calc/certify_crossing.py:1-4 (U95, monotonicity of g_c in p, left Riemann sum plus tail, exact dynamic program), certify_out.json per_threshold (ratio_upper 0.8902 at c = 0.05, grid 6000, h = 6c); the outside check m2-calc/gpt-6.1-sol-check.md (M2-FWER-RECORD.md:23-29: uniform p is the worst case). Methods-paper shortening pass (2026-10-21, DECISION-20261020-joseph-refined-estimator-shorten-delta.md item 2): the article's Appendix A, moved verbatim (labels adapted).
\subsection{Familywise error of the sequential joint-test decisions}
\label{sec:joint5d-fwer}

This derivation supports the familywise statement of \S\ref{sec:joint5d-frozen}.

\emph{Probability model.}  Fix one test, and condition on the data and on the simulated sample.  The calibration pseudo-experiments, taken in their fixed seed order, give exceedance indicators $X_{b}=\mathbf{1}[T_{b}\ge t_{\mathrm{obs}}]$.  Assumption A2 is that, given the data, these are independent Bernoulli draws with common probability $p$, the calibrated tail probability of the observed statistic; this requires the lost pseudo-experiments to have been lost independently of their outcomes, which has not been shown (\S\ref{sec:joint5d-frozen}, which also states the design's two data-informed elements).  Assumption A1 is that, for a true null, $p$ is super-uniform over repetitions of the experiment, $P(p\le u)\le u$, with the nuisance parameters drawn from their priors; the size statement is therefore averaged over the nuisance priors. Let $k_{B}$ be the number of exceedances among the first $B$ draws and $U(k,B)$ the upper end of the two-sided 95\% Clopper--Pearson interval, a one-sided 97.5\% upper bound.

\emph{Crossing event.}  The rule rejects at threshold $c$ only if $U(k_{B},B)<c$ at the number $B\le\jtBcap$ at which sampling stopped.  Whatever rule chose that stopping point, including stops decided with results visible, irregular looks, lost draws, and the budget extensions, the rejection lies in the event $E_{c}=\bigcup_{B\le\jtBcap}\{U(k_{B},B)<c\}$ that the bound crosses below $c$ at some $B$.  Hence $P(\text{reject at }c)\le R(c)\equiv E_{p}[g_{c}(p)]$, with $g_{c}(p)=P(E_{c}\mid p)$.

\emph{Uniform $p$ is the worst case.}  Coupling the Bernoulli draws in $p$ shows that $g_{c}(p)$ does not increase with $p$.  For any super-uniform $p$ its expectation is therefore at most its value for uniform $p$, $R(c)\le\int_{0}^{1}g_{c}(u)\,du$.

\emph{The decision rule, operationally.}  The tests are ordered by increasing $(k+1)/(B+1)$, with ties broken in the fixed listing order (total before shape for each null).  At step $i$ the threshold is $\alpha/(m-i+1)$: the test is rejected if $U<\alpha/(m-i+1)$, it is not rejected if the lower end of its interval exceeds the threshold, and it is undetermined otherwise; a test that is not rejected or undetermined ends the step-down, and the tests after it take the same status.  Equality $U=c$ does not reject.
% X7 (sol Appendix A): M2-AUDIT.md:68 (Holm with determinacy: reject if the upper end < threshold, stop "not rejected" if
% the lower end > threshold, "undetermined" otherwise); a7margins/RECEIPT.md:46-49 (ties by the design's insertion
% order, total before shape); nd-unfolding/s5p_inference.py:127-152.

\emph{Holm reduction without independence.}  Let $m_{0}$ of the $m=10$ hypotheses be true.  The first true null to be rejected is reached after at most $m-m_{0}$ rejections of false nulls, so its threshold is at most $\alpha/m_{0}$.  A familywise error therefore requires some true null to cross below $\alpha/m_{0}$, and by the union bound $\mathrm{FWER}\le m_{0}\,R(\alpha/m_{0})\le\max_{j\le m}j\,R(\alpha/j)$.  The argument uses neither independence between tests nor the order of the Holm steps, so the shared ensemble of a null's total and shape tests, their unequal stopping counts, and the ranking by $(k+1)/(B+1)$ with its ties do not affect it.  The decision $p$ value is the largest over shifted nulls that include the unshifted one, evaluated at the same $B$, so its exceedance count is at least that of the unshifted test; since $U$ increases with $k$, the bound for the unshifted test covers the decision rule.

\emph{Certified constant.}  For each $c=\alpha/j$, $j=1,\dots,10$, let $k^{*}(B)$ be the largest $k$ with $U(k,B)<c$.  For fixed $p$, a dynamic program propagates the distribution of $k_{B}$ from $B=1$ to \jtBcap{}, moving the probability of every state with $k_{B}\le k^{*}(B)$ into an absorbing ``crossed'' state at each $B$, and pooling the counts above $\max_{B}k^{*}(B)$, which can never cross because counts never decrease; the absorbed probability is $g_{c}(p)$ exactly.  This is evaluated on a grid of 6000 points in $[0,6c]$.  Because $g_{c}$ is nonincreasing, a left Riemann sum over-estimates its integral on each grid interval, and with the tail term $(1-6c)\,g_{c}(6c)$ it is a rigorous upper bound on $\int_{0}^{1}g_{c}$, so the integration error is included in the bound. The largest ratio $R(c)/c$, at $j=1$, lies in $[0.8892,0.8902]$; it decreases to $[0.5138,0.5148]$ at $j=10$.  Hence, for $\alpha=0.05$, the value used here, $\mathrm{FWER}\le0.8902\,\alpha<0.9\,\alpha$ under any stopping at or before \jtBcap{} pseudo-experiments.  The bound rests on A1, the calibration's fidelity, which \S\ref{sec:joint5d-conditions} discusses, and on A2.
% X7: m2-calc/certify_crossing.py:1-4, :8-35 (kstar by bisection; absorbing state; overflow pooled; strict U < c).
